\section{Actual symplectic rigidity from adjoint Selmer vanishing}
\label{sec:actual-symplectic-rigidity}

This chapter completes the irreducible-point rigidity argument on the
original arithmetic datum.  The source passes through a substantial
ordinary automorphy-lifting construction.  A more direct route is available:
Thorne's adjoint Selmer theorem requires irreducibility over the cyclotomic
tower, and the ordinary characters of these particular points prove that
condition.  We give the restriction argument, the local crystalline
calculation, and the polarization comparison explicitly.  The resulting
statement supplies the literal \texttt{RigiditySymplecticFour} interface.
It retains the automorphy premise in that interface; automorphy is not
inferred from the existence of a deformation ring.

\subsection{The unchanged datum and the classical inputs}

Fix the original algebraic closure $\Omega/\mathbb Q$, coefficient ring
$O\subset\overline{\mathbb Q}_\ell$, and a datum
$D:\texttt{Raw}(\Omega,O)$ with $D.\texttt{Good}$.  Thus $\ell$ is odd and
$e\ge2$.  Write $G=\operatorname{Gal}(\Omega/\mathbb Q)$,
$\chi:G\to\mathbb Z_\ell^\times$ for the cyclotomic character, and
$c=D.\mathrm{cc}$ for the chosen complex conjugation.  All decomposition
and inertia groups below are the original ones in $G$.  The characters are
\[
 \alpha=\chi\omega^{-1}\xi,\qquad
 \Psi=\chi^e\psi,\qquad \vartheta=\alpha\Psi,
 \qquad (p_0,p_1,p_2,p_3)=(\vartheta,\Psi,\alpha,1).
 \tag{RX.1}
\]
AX6 supplies their exact continuity, finite-order, inertia and separation
properties.  In particular $\vartheta(c)=-1$, and the four inertia weights
on the subgroup where $\psi$ and $\omega$ vanish are
\[
                    (w_0,w_1,w_2,w_3)=(e+1,e,1,0).
 \tag{RX.2}
\]
The fixed alternating matrix and its adjoint operation are
\[
 J=\begin{pmatrix}0&0&0&1\\0&0&1&0\\0&-1&0&0\\-1&0&0&0\end{pmatrix},
 \quad X^\dagger=J^{-1}X^tJ,\quad
 \mathfrak{sp}_4=\{X:X^\dagger=-X\},\quad
 P(X)=\tfrac12(X-X^\dagger).
 \tag{RX.3}
\]

An actual point $r$ is an $O$-algebra specialization of an actual
$\texttt{IsPt}$ lift over a complete Noetherian local $O$-algebra $R$ with
the same residue field.  This is precisely $D.\texttt{PtIn}(A,r)$ for
$A=\overline{\mathbb Q}_\ell$ or
$\overline{\mathbb Q}_\ell[\epsilon]/(\epsilon^2)$.
The local condition is the original \texttt{OnComp} condition, including
its domain and diagonal base point.  In particular AX8 applies to every
such characteristic-zero point.  The source's $\texttt{Irr}(r)$ says that
the matrices $r(g)$ span $M_4(\overline{\mathbb Q}_\ell)$; it implies
absolute irreducibility because every invariant subspace is then invariant
under every matrix.

We use the following named classical inputs, with the specific
normalizations required here.

\begin{itemize}
\item Fontaine's category of crystalline representations is stable under
subquotients, tensor products and duals; its period functors are exact on
admissible objects.  An unramified representation is crystalline, and a
cyclotomic power $\chi^a$ has its de Rham filtration in degree $-a$.
These are the consequences of Proposition~1.5.2 and
\S\S5.1.3--5.1.7 of \cite{RX-Fontaine}.  BX1--BX2 identify these rings,
actions, filtrations and finite cocycles with the literal source models;
BX4 supplies the Bloch--Kato dimension formula.  DX1 supplies finite local
Tate duality and the finite Euler formula; RX2 below spells out the
rational passage for the possibly higher-dimensional representation.
\item Finite abelian class-field theory over $\mathbb Q$, including the
cyclotomic description used in AX3, attaches algebraic Hecke characters
to finite-order characters times integral cyclotomic powers.  We use
geometric Frobenius throughout.  The quadratic character of an imaginary
quadratic extension is trivial after composition with its idele norm.
\item Strong multiplicity one identifies two cuspidal
$\mathrm{GL}_n$ representations whose unramified local parameters agree
outside a finite set; see \cite[\S4]{RX-JS}.
The normalized Galois representation of a regular algebraic polarized
cuspidal representation is the one in
\cite[Theorem~2.1.1]{RX-BLGGT}.  Its geometric Frobenius eigenvalues are
$q^{(n-1)/2}$ times the unitary Satake parameters.  We use Chebotarev
density and the characteristic-zero character criterion for semisimple
representations to identify two such attached systems.
\item The substantive global vanishing input is
\cite[Theorem~6.1]{RX-Thorne}: if $F$ is CM, $n\ge2$, and
$(\pi,\lambda)$ is regular algebraic, cuspidal and polarized, then
irreducibility of $r_{\pi,\iota}$ on
$G_{F(\mu_{\ell^\infty})}$ implies
$H^1_f(F^+,\operatorname{ad}r_{\pi,\iota})=0$.
Here the adjoint is the polarized, or unitary, extension to $G_{F^+}$.
It will be identified by an explicit formula below.  The theorem has no
residual irreducibility or residual image-size premise.
\end{itemize}

The standard results are leaves of this reconstruction.  The class-field
and cyclotomic facts in the second item are the already matched AX3
inputs.  The enlarged
source packages \texttt{UnramifiedInputs} and
\texttt{AdjointLiftingTheorem} are not assumed in the argument below and
are not claimed as consequences of it.

\subsection{Actual coefficients and ordinary finite cocycles}

% BEGIN RX1
\begin{proposition}[RX1: finite coefficients and the actual tangent cocycle]
\label{prop:rx-coefficients}
Every actual $\overline{\mathbb Q}_\ell$ point is continuous and defined
over one finite extension of $\mathbb Q_\ell$.  The same holds for both
coordinates of every actual dual-number point.  If $r'$ is such a point
and its first coordinate is $r$, write
\[
 r'(g)=r(g)+\epsilon\dot r(g),\qquad
 z(g)=\dot r(g)r(g)^{-1}.
 \tag{RX.4}
\]
Then $z$ is a continuous adjoint cocycle, and every $z(g)$ belongs to the
actual $\mathfrak{sp}_4$ in (RX.3).  Moreover $r$ is unramified away from
a finite set of rational primes.
\end{proposition}

\begin{proof}
We first prove the coefficient assertion for the algebra map defining the
point.  Let $f:R\to\overline{\mathbb Q}_\ell$ be an $O$-algebra map.
For $a\in\mathfrak m_R$ one has $|f(a)|<1$.  Otherwise the inverse of
the nonzero algebraic number $f(a)$ is integral over $\mathbb Z_\ell$.
Its monic polynomial, multiplied by a suitable power of $f(a)$, gives
\[
              1+b_1f(a)+\cdots+b_df(a)^d=0,
                   \qquad b_i\in\mathbb Z_\ell.
 \tag{RX.5}
\]
The corresponding element of $R$ lies in $1+\mathfrak m_R$ and is a
unit, which cannot map to zero.  Every element of $R$ differs from an
element of $O$ by an element of $\mathfrak m_R$, so $|f(R)|\le1$.
Choose finitely many generators of $\mathfrak m_R$ and a real number
$0<a<1$ bounding their image norms.  The ultrametric inequality gives
\[
                 |f(x)|\le a^m\quad(x\in\mathfrak m_R^m).
 \tag{RX.6}
\]
Thus $f$ is continuous.  Lifting generators of
$\mathfrak m_R/(\mathfrak m_R^2+\mathfrak m_O R)$ and successive
adic approximation gives a surjection $O[[t_1,\ldots,t_d]]\to R$.
Choose a finite extension $E/\mathbb Q_\ell$ containing
$\operatorname{Frac}(O)$ and the finitely many $f(t_i)$.
Polynomial approximations converge after applying $f$ by (RX.6).
Since $E$ is complete and is closed in $\overline{\mathbb Q}_\ell$,
they show $f(R)\subset O_E$.

Now let $F=f+\epsilon\delta:R\to
\overline{\mathbb Q}_\ell[\epsilon]/(\epsilon^2)$ be an $O$-algebra
map.  The function $\delta$ is an $O$-linear $f$-derivation.
The ideal $I=\ker f$ has finitely many generators $a_j$.
The $O$-submodule $f(R)\subset O_E$ is finite; choose generators
$f(b_i)$.  Every $r\in R$ can therefore be written
\[
             r=\sum_i o_i b_i+\sum_j u_j a_j,
                      \qquad o_i\in O,\quad u_j\in R.
\]
Consequently
\[
             \delta(r)=\sum_i o_i\delta(b_i)
                         +\sum_j f(u_j)\delta(a_j).
 \tag{RX.7}
\]
Adjoining the finitely many displayed derivative values to $E$ gives one
finite coefficient field $E'$ containing all of $F(R)$, and (RX.7)
gives a uniform bound $|\delta(R)|\le C$.
Leibniz's rule and (RX.6) then give
$|\delta(x)|\le Ca^m$ on $\mathfrak m_R^{m+1}$.
The coefficients multiplying products of generators cause no problem:
their derivative term has one additional factor bounded by $a$.
Thus $\delta$ is also adically continuous.  This proves finite
coefficient capture and continuity for the original maps themselves;
it does not use any assertion that a compact subset of
$\overline{\mathbb Q}_\ell$ lies in a finite extension.

The actual lift has continuous entries modulo every power of
$\mathfrak m_R$.  Applying the preceding estimates proves continuity
of $r$ and $\dot r$ in their finite coefficient fields.  Multiplication
of the dual-number matrices gives
\[
 z(gh)=z(g)+r(g)z(h)r(g)^{-1}.
 \tag{RX.8}
\]
The multiplier is fixed by the $O$-algebra map, so the coefficient of
$\epsilon$ in
$r'(g)^tJr'(g)=\vartheta(g)J$ is
$z(g)^tJ+Jz(g)=0$.  This is exactly $z(g)^\dagger=-z(g)$.

Finally the residual representation has open kernel and finite image
by AX6.  Its finite fixed field is ramified at finitely many primes.
Outside these primes and $D.T$, its entire inertia group lies in
$\ker\bar r$, and the original \texttt{UnrT} equation says that the
actual lift, hence $r$, is trivial on that inertia group.
\end{proof}
% END RX1

% BEGIN RX2
\begin{proposition}[RX2: the ordinary crystalline calculation and descent]
\label{prop:rx-ordinary-crystalline}
Let $L/\mathbb Q_\ell$ and $E/\mathbb Q_\ell$ be finite, and let
$A=E$ or $E[\epsilon]/(\epsilon^2)$.  Suppose a continuous upper
triangular $A$-representation of $G_L$ has diagonal characters
$\kappa_1,\ldots,\kappa_n$ such that
\[
 \kappa_i|_{I_L}=\chi^{w_i},\qquad
 w_1>\cdots>w_n,\qquad
 \bar\kappa_i\ne\bar\kappa_j\chi\quad(i<j).
 \tag{RX.9}
\]
Here the bar means reduction modulo $\epsilon$ when $A$ has dual
numbers, and is the identity otherwise.  Its underlying
$\mathbb Q_\ell$ representation is crystalline.
In particular a dual-number self-extension of a representation that
satisfies (RX.9) has finite adjoint cocycle.  If such a cocycle becomes
finite on an open subgroup, it is finite on the original local group.
\end{proposition}

\begin{proof}
Each rank-one $A$-character $\kappa_i\chi^{-w_i}$ is unramified.
As an $E$-representation it may be a nontrivial unramified extension
of a character by itself, but is still crystalline: its restriction
to inertia is trivial, and Fontaine's unramified descent
\cite[\S5.1.5]{RX-Fontaine} gives crystalline admissibility.
Twisting by $\chi^{w_i}$ preserves it.

Induct on $n$.  The first $n-1$ coordinates give a crystalline
$A$-module $V$ by induction.  Twisting by $\kappa_n^{-1}$ gives a
crystalline underlying $E$-representation
$W=V\otimes_A A(\kappa_n^{-1})$.
Indeed this tensor product is a quotient of
$V\otimes_E A(\kappa_n^{-1})$, so the assertion follows from
tensor stability and subquotient stability of crystalline representations.
It has an $E$-character filtration with factors
$\bar\kappa_i/\bar\kappa_n$ for $i<n$, each occurring twice when
$A=E[\epsilon]$ and once when $A=E$.
Their inertia exponents are positive.  Exactness and strictness of
$D_{\mathrm{dR}}$ therefore give
\[
 \operatorname{Fil}^0D_{\mathrm{dR}}(W)=0,
 \qquad \dim_E t_W=[L:\mathbb Q_\ell]\dim_E W.
 \tag{RX.10}
\]
The constituents of $W^\vee(1)$ are all nontrivial by (RX.9).
Successive left exactness of invariants gives $H^0(L,W^\vee(1))=0$.
We record the needed higher-dimensional rational Euler passage.
Choose a stable $\mathbb Z_\ell$-lattice $T\subset W$; it exists
because the continuous image of the compact group is bounded, so
the sum of the translates of a lattice is contained in a fixed
larger lattice.  The finite-cochain limit and finite-generation
argument of BX3 applies in every degree at issue.  It gives
$H^i(L,T)\otimes\mathbb Q_\ell=H^i(L,W)$ and finite generation
of $H^i(L,T)$ for $i=0,1,2$, with $H^3(L,T)=0$.
The multiplication-by-$\ell^m$ exact sequence gives
\[
 0\longrightarrow H^i(L,T)/\ell^m
 \longrightarrow H^i(L,T/\ell^mT)
 \longrightarrow H^{i+1}(L,T)[\ell^m]\longrightarrow0.
\]
Its outside torsion groups have bounded size independently of $m$.
Thus the base-$\ell$ logarithm of the middle cardinality is
$m\operatorname{rank}_{\mathbb Z_\ell}H^i(L,T)+O(1)$.
Divide the finite Euler identity of DX1 by $m$ and let $m$ increase.
This proves the rational Euler identity with factor
$[L:\mathbb Q_\ell]\dim_{\mathbb Q_\ell}W$.
The degree-$(2,0)$ finite Tate pairings similarly identify
$H^2(L,T)$ with the compact dual of $H^0(L,W^\vee(1)/T^\vee(1))$.
If $H^0(L,W^\vee(1))=0$, the latter invariant group embeds into
the finite torsion of $H^1(L,T^\vee(1))$, by the lattice sequence.
Hence $H^2(L,W)=0$.
After dividing all dimensions by $[E:\mathbb Q_\ell]$, these
facts and BX4 yield
\[
 \begin{split}
 H^2(L,W)&=0,\\
 \dim_E H^1(L,W)&=\dim_E H^0(L,W)
                              +[L:\mathbb Q_\ell]\dim_E W\\
                    &=\dim_E H^1_f(L,W).
 \end{split}\tag{RX.11}
\]
These are finite-dimensional spaces; hence $H^1_f(L,W)=H^1(L,W)$.
The upper-right column of the representation, divided by
$\kappa_n$, is a cocycle in $W$.  Its image in
$\operatorname{Hom}_E(A(\kappa_n),V)$ is the class of the
underlying $E$-extension.  A $B_{\mathrm{cris}}$ primitive in $W$
maps to a primitive in this Hom space.  After tensoring with
$B_{\mathrm{cris}}$ the extension therefore splits.  Since its
subrepresentation and quotient are crystalline, their comparison
isomorphisms give one for the whole extension.  This completes the
induction.  In particular it includes all dual-number unramified
diagonal variations, rather than assuming that those variations vanish.

For a crystalline self-extension
$0\to V\to V'\to V\to0$, exactness of $D_{\mathrm{cris}}$
and the comparison isomorphisms give a $G_L$-equivariant splitting
after tensoring with $B_{\mathrm{cris}}$.
One can choose the splitting over the semisimple coefficient algebra
$L_0\otimes_{\mathbb Q_\ell}E$ before extending scalars.
Writing it in a chosen vector-space splitting says precisely that
the adjoint cocycle has a $B_{\mathrm{cris}}$ primitive.  BX1--BX2
identify this with the source's \texttt{finiteCocycles}, including
its specified tensor order and diagonal action.

Here is the descent argument at the cochain level.  Shrinking an open
subgroup to its normal core, suppose $H\lhd G_L$ is open and
$z(h)=hx-x$ for $h\in H$, with
$x\in V\otimes_{\mathbb Q_\ell}B_{\mathrm{cris}}$.
The cocycle $F(g)=z(g)-(gx-x)$ vanishes on $H$.  The cocycle identity
shows that it factors through $Q=G_L/H$ and takes values in the
$H$-invariants.  Put $S=\sum_{q\in Q}F(q)$.  Then
\[
     gS-S=-|Q|F(g),\qquad
     z(g)=g\bigl(x-|Q|^{-1}S\bigr)
                         -\bigl(x-|Q|^{-1}S\bigr).
 \tag{RX.12}
\]
The division is valid in characteristic zero, including when
$\ell\mid|Q|$.  This proves descent without an unramified-extension
restriction.  Equivariant coefficient maps and projections carry
the displayed primitive to the corresponding primitive, so the same
argument also proves the coefficient transports used below.
\end{proof}
% END RX2

% BEGIN RX3
\begin{proposition}[RX3: the full actual \texttt{tanBK} conclusion]
\label{prop:rx-tanbk}
Let $D$ be good, let $r$ be an actual point, and let $r'$ be any
actual dual-number point reducing to $r$.  The cocycle in (RX.4)
satisfies, on the original $G$,
\[
 z\in Z^1(G,\operatorname{Ad}r),\qquad
 \texttt{Af.IsBK}(\operatorname{Ad}r,z),\qquad
                   z(g)\in\mathfrak{sp}_4\quad(g\in G).
 \tag{RX.13}
\]
Thus the source's \texttt{tanBK} output is supplied with every
original local condition, coefficient map and placement retained.
\end{proposition}

\begin{proof}
RX1 gives the first and third assertions.  Enlarge its single finite
coefficient field if needed for the finitely many changes of basis
below; the original embeddings into $\overline{\mathbb Q}_\ell$ remain
fixed.

At $\ell$, the actual \texttt{OnComp} condition implies the
unnormalized ordinary polynomial equations.  The four prescribed
eigenvalues $p_i(s_0)$ are distinct by AX6, so their differences are
units after passage to $E[\epsilon]$.
The spectral projectors
\[
        \prod_{j\ne i}
           \frac{r'(s_0)-p_j(s_0)}{p_i(s_0)-p_j(s_0)}
 \tag{RX.14}
\]
decompose the free module into rank-one summands.  Their partial sums
are stable by the ordinary equations.  A basis of these summands
therefore makes all $r'(g)$ upper triangular, with continuous
diagonal characters $\kappa_i$ and
$\kappa_i|_{I_\ell}=p_i|_{I_\ell}$ exactly.
Reduction gives an ordinary flag for $r$.  Its ordered characters
are those of AX8: Jordan--H\"older uniqueness identifies the sets of
characters, and evaluation at the distinct prescribed eigenvalues
of $s_0$ identifies the indices.  No permutation of the ordinary
weights has been introduced.

Put $\sigma=D.\mathrm{sg}(\ell)$.
The finite character $\tau=(\psi,\omega)$ of the original local
group has an unramified companion $\tau_{\mathrm{ur}}$ taking
$\sigma$ to $\tau(\sigma)$.  Indeed
$G_\ell/I_\ell=\widehat{\mathbb Z}$ and a chosen Frobenius is its
topological generator, so a homomorphism to the finite abelian image
is determined by this value.  Therefore
\[
 U=\ker(\tau\tau_{\mathrm{ur}}^{-1})
 \tag{RX.15}
\]
is open and normal, contains $\sigma$, and kills both $\psi$ and
$\omega$ on $U\cap I_\ell$.
The inertia characters of $r'|_U$ are consequently the exact powers
$\chi^{e+1},\chi^e,\chi,1$.
For every $i<j$, AX8, applied to this very subgroup $U$, says that
$\bar\kappa_i|_U\ne\bar\kappa_j|_U\chi$.
The presence of $\sigma$ in $U$ is essential in this use of AX8;
distinct integer weights alone would not exclude the two adjacent
weight differences equal to one.

The actual decomposition group, its henselization and completion
identify $U$ with the absolute Galois group of a finite local
extension.  The inertia and cyclotomic identifications are those
already established in DX1 and transported to the literal periods
in BX1--BX2.  RX2 makes the underlying $r'$ and $r$ crystalline on
$U$, so $z|_U$ is finite.  Formula (RX.12) then returns the finite
cocycle to the entire original decomposition group.  The period
primitive is transported through the original graph subgroup
$\texttt{decJ}$, giving the exact \texttt{IsBK.above} clause.

For a prime $q\ne\ell$ in $D.S_b$, the original condition kills
$I_q\cap\ker\psi$.  Away from $D.T$, it kills
$I_q\cap\ker\bar r$.  Both are finite-index normal subgroups of
inertia, and $r'$ is the identity on the relevant subgroup.
Thus $z$ vanishes there and factors, on inertia, through a finite
quotient.  Averaging this finite quotient as in (RX.12), now with
coefficients in $M_4(E)$, shows that $z|_{I_q}$ is a coboundary.
For every rational finite place one may choose the original selected
valuation if $q\in D.S_b$, and any original valuation above it
otherwise.  Conjugating a placement transports a cocycle and its
primitive by the cocycle identity.  This is precisely the
existential original-valuation clause \texttt{IsBK.away}, at every
place, and proves (RX.13).
\end{proof}
% END RX3

\subsection{Cyclotomic irreducibility and an admissible CM field}

% BEGIN RX4
\begin{proposition}[RX4: irreducibility over the required infinite tower]
\label{prop:rx-cyclotomic-irreducibility}
If $D$ is good and $r$ is an irreducible actual point, then
$r|_{\ker\chi}$ is irreducible.  There is an imaginary quadratic
field $K\subset\Omega$ such that
\[
       r|_{G_{K(\mu_{\ell^\infty})}}\text{ is irreducible}.
 \tag{RX.16}
\]
In particular $r|_{G_K}$ is irreducible.  The field $K$ is chosen
inside the original $\Omega$; it does not replace the original
arithmetic datum or its coefficient field.
\end{proposition}

\begin{proof}
We need the following elementary profinite Clifford argument.
Let $\Gamma$ act continuously and irreducibly on a finite-dimensional
space over $\overline{\mathbb Q}_\ell$, let $H\lhd\Gamma$ be closed,
and suppose $\Gamma/H$ is procyclic.  If the restriction to $H$ is
reducible, there is a nontrivial finite-order continuous character
$\eta$ of $\Gamma/H$ such that
\[
                          r\otimes\eta\simeq r.
 \tag{RX.17}
\]
Here is a proof that also covers an infinite subgroup $H$.
Choose an $H$-invariant nonzero subspace of minimal dimension.  It
is simple.  Its $\Gamma$-translates are simple $H$-submodules, and
their sum is a nonzero $\Gamma$-invariant subspace, hence the whole
space.  A sum of simple modules is semisimple: successively add a
simple summand only when it is not already contained in the sum;
its intersection with that sum is then zero.  Finite dimension
ends the procedure.  Thus the restriction has finitely many
isotypic components, permuted transitively by $\Gamma$.

Each isotypic component is a closed vector subspace.  Its stabilizer
in $\Gamma$ is closed, since the condition that a continuous matrix
map preserve a given finite-dimensional subspace is closed.  Its
index is finite, so its complement is a finite union of closed
cosets and the stabilizer is open.  The action on the components
is consequently a continuous finite quotient of the procyclic
group.  If the number of components is $d>1$, a generator cycles
them.  Multiplication by $1,\zeta_d,\ldots,\zeta_d^{d-1}$ on
successive components gives an intertwiner for a nontrivial
character of this cyclic quotient, proving (RX.17).

It remains to exclude multiplicity without a component permutation.
In that case the restriction is $W\otimes M$, with $W$ simple for
$H$ and $H$ acting trivially on the multiplicity space $M$.
Choose $t\in\Gamma$ whose image topologically generates
$\Gamma/H$.  The conjugate of $W$ by $t$ is isomorphic to $W$;
choose an invertible intertwiner $A$ for this conjugation.
Schur's lemma gives $\operatorname{End}_H(W)=
\overline{\mathbb Q}_\ell$, so the action of $t$ is $A\otimes B$
for some invertible $B$ on $M$.
An eigenline of $B$ makes $W\otimes L$ invariant under $H$ and $t$.
The stabilizer of this subspace is closed and contains $H$ and a
dense generator modulo $H$, hence is all of $\Gamma$.
Irreducibility forces $\dim M=1$.  This proves the assertion, with
no finite-group Clifford theorem silently applied to an infinite
group.

Apply it to $\Gamma=G$ and $H=\ker\chi$.
For odd $\ell$,
$\mathbb Z_\ell^\times\simeq C_{\ell-1}\times\mathbb Z_\ell$
is procyclic, and the global cyclotomic character is surjective.
Suppose that (RX.17) holds.  At the original group $G_\ell$,
AX8 gives the ordinary characters $\kappa_0,\ldots,\kappa_3$.
The two isomorphic representations have the same semisimplification,
so their constituents satisfy
\[
               \{\kappa_i\eta:0\le i<4\}
                       =\{\kappa_i:0\le i<4\}.
 \tag{RX.18}
\]
Restrict further to an open subgroup of inertia killing $\eta$,
$\psi$ and $\omega$.  The four characters there are the distinct
powers (RX.2); the image of $\chi$ is infinite on every open
inertia subgroup.  Equality of two powers therefore forces equality
of their integer exponents.  Thus the permutation in (RX.18)
fixes every index, and $\eta|_{G_\ell}=1$.

The image of $I_\ell$ under $\chi$ is all of
$\mathbb Z_\ell^\times$.  One elementary verification is that
$\Phi_{\ell^m}(1+T)$ is Eisenstein at $\ell$; hence the local
cyclotomic extension is totally ramified of degree
$\varphi(\ell^m)$.  The original global inertia therefore surjects
onto every finite cyclotomic Galois group.  Compactness and the
inverse-limit description give the claimed full image.
A character of $G/H$ trivial on $I_\ell$ is consequently trivial.
This contradicts (RX.17), proving irreducibility on $H$.

By RX1 choose an odd prime $q\ne\ell$ at which $r$ is unramified,
and take $K=\mathbb Q(\sqrt{-q})$ inside $\Omega$.
Its quadratic character $\eta_K$ is nontrivial on $I_q$, since
$q$ has odd valuation in $-q$ and ramifies in $K$.
The cyclotomic character is unramified at $q$, so $I_q\subset H$.
Hence $H\cap G_K$ has index two in $H$.
If the restriction of the irreducible $r|_H$ to this subgroup were
reducible, the same Clifford argument for this cyclic quotient of
order two would give
$r|_H\simeq r|_H\otimes\eta_K|_H$.
On $I_q$ the first representation is the identity, while the second
contains the nontrivial scalar $-1$.  This is impossible.
Finally $H\cap G_K=G_{K(\mu_{\ell^\infty})}$, which proves
(RX.16).
\end{proof}
% END RX4

\subsection{The exact automorphic and adjoint comparison}

% BEGIN RX5
\begin{proposition}[RX5: the polarized adjoint of the actual automorphic point]
\label{prop:rx-adjoint-vanishing}
Suppose $D$ is good, $r$ is an irreducible actual point, and
$\texttt{IsAut4Alt}(\iota,r)$ holds with its literal source meaning
$\texttt{IsAutQA}$.  Let $K$ be as in RX4, and let $\eta_K$ be
its quadratic character.  On $M_4(\overline{\mathbb Q}_\ell)$ define
\[
 \mathcal R(g)X
   =r(g)P(X)r(g)^{-1}
       +\eta_K(g)r(g)(X-P(X))r(g)^{-1}.
 \tag{RX.19}
\]
This is the polarized adjoint representation of the automorphic
representation supplied by \texttt{IsAut4Alt}, with the same
original $G$ and $r|_{G_K}$.  Its actual Bloch--Kato Selmer group
over $\mathbb Q$ vanishes.
\end{proposition}

\begin{proof}
RX1 supplies continuity, and RX4 supplies irreducibility of
$r|_{G_K}$, so the universal imaginary-quadratic clause of
$\texttt{IsAutQA}$ applies to this $K\subset\Omega$.
It gives a regular algebraic cuspidal representation $\pi/K$
whose \texttt{AttC} system is $r|_{G_K}$.
The additional \texttt{WitA} assertion in the literal definition
is not needed in this argument.

We must first verify polarization with the normalization of that
system.  Since $r(g)^tJr(g)=\vartheta(g)J$, the representation
is self-dual up to $\vartheta$; conjugation by $c$ also gives an
isomorphism from $r|_{G_K}$ to its conjugate.
At an unramified place $v$ of norm $Q_v$, write the unitary
Satake parameters of $\pi_v$ as $a_{v,i}$.
The definition of \texttt{AttC} uses geometric Frobenius $F_v$
and gives the eigenvalues $Q_v^{3/2}a_{v,i}$.
Because $\chi(F_v)=Q_v^{-1}$, the preceding self-duality and
conjugation identities yield the exact multiset relation
\[
 \{a_{v^c,i}\}_i
       =\{\iota(\vartheta(F_v))Q_v^{-3}a_{v,i}^{-1}\}_i.
 \tag{RX.20}
\]
The character $\vartheta\chi^{-(e+1)}$ has finite order and open
kernel.  Class-field theory therefore supplies an algebraic
Hecke character $\lambda_0$ of $\mathbb Q$ whose pullback to $K$
has $\ell$-adic character
\[
                    (\vartheta\chi^3)|_{G_K}.
 \tag{RX.21}
\]
Concretely this is the finite Dirichlet character supplied by the
cyclotomic description of the finite-order part, multiplied by
the integral norm power fixed by its geometric Frobenius values.
Thus its local factor in (RX.20) is exactly
$\iota(\vartheta(F_v))Q_v^{-3}$.
Strong multiplicity one, in the form of
\cite[Theorem~4.2]{RX-JS}, gives
\[
       \pi^c\simeq\pi^\vee\otimes
                    (\lambda_0\circ N_{K/\mathbb Q}\circ\det).
 \tag{RX.22}
\]
If needed, multiply $\lambda_0$ by the quadratic Hecke character
of $K/\mathbb Q$ so that its value at $-1$ at the real place is
$(-1)^4=1$.  That character has value $-1$ there and is trivial
after the idele norm from $K$.  Hence neither (RX.21) nor (RX.22)
changes.  This is the exact sign correction in the source's
\texttt{heckeChar\_sign\_fixed}; no equality between a Hecke
infinity sign and an uncorrected cyclotomic exponent is assumed.
The resulting pair $(\pi,\lambda)$ is polarized in the precise
sense of \cite[\S2.1]{RX-BLGGT}.

The normalized Galois representation of this pair has the same
geometric Frobenius characteristic polynomials as $r|_{G_K}$
outside a finite set, by (RX.20)--(RX.21) and
\cite[Theorem~2.1.1]{RX-BLGGT}.
Both representations are continuous and semisimple, so Chebotarev
\cite[Theorem~8.31]{RX-Milne}
and the character criterion identify them.  To see that the use
of $\overline{\mathbb Q}_\ell$ hides no change of topology,
choose finite coefficient fields for both representations, take
their compositum, and compare their continuous traces there.
The finite-dimensional character criterion can be seen by taking the
algebra generated by the two actions in
$\operatorname{End}(V\oplus V')$.  Semisimplicity and faithfulness
make this a semisimple algebra.  After scalar extension its simple
matrix factors have linearly independent trace characters, detected
by their central idempotents.  Equality of traces therefore forces
equality of the multiplicities of every simple module, and hence
an isomorphism.  In particular RX4 proves exactly the
cyclotomic irreducibility hypothesis in Thorne's theorem for this
attached representation.

It remains to identify its extension of the adjoint to $G$.
Put
\[
                            B=Jr(c)^{-1}.
 \tag{RX.23}
\]
Since $c^2=1$ and $\vartheta(c)=-1$, the matrix $B$ is symmetric.
For $s\in G_K$, direct multiplication gives
\[
         r(s)^t B r(csc^{-1})=\vartheta(s)B.
 \tag{RX.24}
\]
Thus the polarized adjoint acts by ordinary conjugation on $G_K$
and by
\[
                           X\longmapsto-B^{-1}X^tB
 \tag{RX.25}
\]
at $c$.  These are the usual Lie-algebra formulas for the polarized
extension to $\mathcal G_4$.  Any symmetric polarization pairing
with the same multiplier on $G_K$ differs from $B$ by a scalar:
the ratio gives an endomorphism of the irreducible $r|_{G_K}$.
Thus this formula agrees with the pairing used for the normalized
automorphic representation, and is independent of its scalar
normalization.

For completeness (RX.19) gives that same representation without
choosing a basis for $B$.  Conjugation by a symplectic similitude
commutes with $\dagger$, hence with the idempotent $P$.
The two summands $\operatorname{im}P$ and $\ker P$ are invariant;
(RX.19) is their ordinary and quadratic-twisted action, so is a
representation.  On $G_K$ it is ordinary conjugation.  At $c$,
\[
 P(X)-(X-P(X))=-X^\dagger
\]
makes (RX.19) equal to (RX.25), using (RX.23).
Since $G_K$ and $c$ generate $G$, the representations agree.
Now apply \cite[Theorem~6.1]{RX-Thorne}, with $F=K$, $F^+=\mathbb Q$
and $n=4$.  It gives $H^1_f(\mathbb Q,\mathcal R)=0$.
The local condition away from $\ell$ is restriction to inertia,
and the condition at $\ell$ is a $B_{\mathrm{cris}}$ primitive.
These are exactly \texttt{Af.IsBK}, by BX1--BX2 and the original
decomposition-group comparison.  At the real place positive
cohomology vanishes by averaging, since the coefficients have
characteristic zero.  There is no additional omitted archimedean
condition.

We also justify the literal algebraic-closure coefficient scope.
The action $\mathcal R$ is defined over a finite field $E_0$ by
RX1.  For any continuous $\mathcal R$-cocycle $a$, use the countable
cofinal family of finite fields constructed in RV4, enlarged to
contain $E_0$.  The sets
$H_i=\{g:a(g)\in M_4(E_i)\}$ are closed subgroups: the cocycle
and inverse formulas preserve these matrices because all action
coefficients lie in $E_0$.  They cover the compact group $G$.
Baire's theorem makes one $H_i$ open; adjoining the finitely many
matrix entries of $a$ on representatives of $G/H_i$ then captures
all its values in one finite extension $E$.
For a local primitive with coefficients in the algebraic closure,
its finitely many coefficient entries lie in some finite $E'/E$.
An $E$-linear retraction $E'\to E$ fixing $1$ commutes with the
action and returns that primitive to $M_4(E)$; at $\ell$ tensor
the retraction with $1$ on $B_{\mathrm{cris}}$.
This is done separately at each place and does not require a
single field containing infinitely many chosen local witnesses.
Thus every cocycle in the literal \texttt{Af.IsBK} condition is
already in the usual finite-coefficient Selmer group after one
finite extension.  Thorne's vanishing gives its actual coboundary
there and hence in the original algebraic closure.
\end{proof}
% END RX5

% BEGIN RX6
\begin{theorem}[RX6: literal actual irreducible rigidity]
\label{thm:rx-literal-rigidity}
For every good original datum $D$, every embedding
$\iota:\overline{\mathbb Q}_\ell\simeq\mathbb C$, and every
irreducible actual point $r$,
\[
          \texttt{IsAut4Alt}(\iota,r)
                       \quad\Longrightarrow\quad D.\texttt{RigidAt}(r).
 \tag{RX.26}
\]
Consequently the universally quantified source proposition
\texttt{RigiditySymplecticFour} is supplied.
\end{theorem}

\begin{proof}
Let $r'$ be any actual dual-number point reducing entrywise to $r$.
RX3 supplies the original cocycle $z$ and all its finite local
conditions.  It takes values in $\operatorname{im}P$, and
\[
 P\circ\operatorname{Ad}_r(g)
             =\mathcal R(g)\circ P,\qquad Pz=z.
 \tag{RX.27}
\]
Apply $P$ to every cochain identity and to every local primitive.
For the prime-$\ell$ condition use $P\otimes1$ on the actual
$B_{\mathrm{cris}}$ tensor module.  Formula (RX.27) proves directly
that the unchanged function $z$ is a cocycle for $\mathcal R$ and
satisfies its actual \texttt{IsBK} conditions.

RX5 gives a matrix $X$ with $z(g)=\mathcal R(g)X-X$.
Apply $P$ once more and multiply on the right by $r(g)$:
\[
    \dot r(g)=r(g)P(X)-P(X)r(g)
                   =Y r(g)-r(g)Y,\qquad Y=-P(X)\in\mathfrak{sp}_4.
 \tag{RX.28}
\]
This proves exactly the existential matrix and the entrywise
identity required by $D.\texttt{RigidAt}(r)$, for every allowed
$r'$; no normalization or choice of its local frames was added.

The literal \texttt{RigiditySymplecticFour} statement assumes
$\texttt{Ctx2}(D)\lor\texttt{Ctx3}(D)$, a positive level $N$
killing the three finite characters, and then the displayed
actual-point and automorphy premises.  Either context includes
$D.\texttt{Good}$.  Applying (RX.26) therefore proves all its
quantifiers, with the same original $D$, $\Omega$, $O$, $\iota$
and $r$.  The unused extra level hypotheses need not be strengthened.
\end{proof}
% END RX6

% BEGIN RX7
\begin{corollary}[RX7: installation on the same selected datum]
\label{cor:rx-same-datum}
Fix an original host $h$, a threshold $t$ with $h.\texttt{Big}(t)$,
an odd prime $\ell$, the original coefficient ring $O$, an integer
$0<e<t$, and $\psi\in h.X(O,e)$.  Let $D$ be the original datum
with $D.\texttt{Good}$ and the literal
$\texttt{OfCase}(D,h,e,\psi)$ in either paired-count case.
The literal \texttt{SymplecticFourImageDense} premise supplies an
embedding $\iota$ and
$D.\texttt{AllIrr}(\texttt{IsAut4Alt}(\iota,\cdot))$ on that
same datum.  It follows that
\[
                          D.\texttt{AllIrr}(D.\texttt{RigidAt}).
 \tag{RX.29}
\]
In the already reconstructed paired-count chain, the independent
actual-rigidity premise can therefore be removed from the cut:
the remaining automorphic requirements are the same-datum image
density supplier and the exact finite catalogue with its trace
nonvanishing guard.
\end{corollary}

\begin{proof}
Choose the one embedding supplied by image density.  For every
$r$, its actual-point and irreducibility assumptions give the
corresponding \texttt{IsAut4Alt} assertion at this embedding.
RX6 proves $D.\texttt{RigidAt}(r)$ pointwise, which is exactly
the quantifier order in \texttt{AllIrr}.
This argument changes neither the selected coefficient ring nor
any local subgroup.  AX10, the actual model constructor MX10,
the model comparison PX1--PX10 and SP10 can therefore consume the
result on their existing common $D$.  Their coefficient-extension
and return steps, the separate \texttt{Fin0}/\texttt{NoEnd}
consequences of image density, and the catalogue's guarded trace
quantifiers retain precisely their previous scope.
No assertion is made here that image density or that catalogue
holds without its own supplier, or that either count or the full
RH/GRH argument is thereby unconditional.
\end{proof}
% END RX7

\subsection{Source boundary and bibliography}

The source chain begins at \texttt{RigSp4.\allowbreak rigSp4\_alt}.
Its auxiliary CM-field theorem obtains adjoint vanishing from
\texttt{UnramifiedInputs} and
\texttt{Adjoint\allowbreak Lifting\allowbreak Theorem}.
The exact declaration and its proof range are recorded in the source ledger.
RX4--RX5 replace that intermediate supplier by the stated
cyclotomic irreducibility proof and Thorne's published theorem.
RX1--RX3 reconstruct the actual \texttt{CopNormLt.\allowbreak coPt},
\texttt{TbTriDiag.\allowbreak tanBK}, ordinary crystalline induction and
finite-index primitive descent.  RX6 uses exactly the projection
and sign in \texttt{RigSp4.\allowbreak rg\_alg}.  Source anchors distinguish
the literal endpoints from the intentionally bypassed larger
interfaces.  In particular this chapter does not promote the
whole Hida automorphy-lifting, unramified tower, or modularity
packages.

\begin{thebibliography}{9}
\bibitem{RX-Fontaine}
J.-M. Fontaine, \emph{Repr\'esentations $p$-adiques semi-stables},
Ast\'erisque \textbf{223} (1994), 113--184,
Proposition~1.5.2 (printed p.~124; PDF index 12),
\S\S5.1.3--5.1.7 (printed pp.~155--158; PDF indices 43--46).
\url{https://www.numdam.org/item/AST_1994__223__113_0.pdf}.

\bibitem{RX-JS}
H. Jacquet and J. A. Shalika,
\emph{On Euler Products and the Classification of Automorphic Forms II},
Amer. J. Math. \textbf{103} (1981), 777--815, Theorem~4.2,
printed pp.~806--807 (PDF indices 30--31); its linear independence
statement implies the strong multiplicity-one form used here.
\url{https://www.math.columbia.edu/~hj/On%20Euler%20products%20II.pdf}.

\bibitem{RX-BLGGT}
T. Barnet-Lamb, T. Gee, D. Geraghty and R. Taylor,
\emph{Potential automorphy and change of weight},
Ann. of Math. \textbf{179} (2014), 501--609,
\S2.1 and Theorem~2.1.1, printed pp.~535--538
(published PDF indices 34--37).
\url{https://annals.math.princeton.edu/wp-content/uploads/annals-v179-n2-p03-p.pdf}.

\bibitem{RX-Thorne}
J. A. Thorne, \emph{On the vanishing of adjoint Bloch--Kato Selmer
groups of irreducible automorphic Galois representations},
Pure Appl. Math. Q. \textbf{18} (2022), no.~5,
Theorem~6.1; author preprint printed p.~24 (PDF index 23).
The definition of the polarized adjoint is recalled on pp.~3 and 25.
\url{https://arxiv.org/pdf/2207.04925}.
DOI: \url{https://doi.org/10.4310/PAMQ.2022.v18.n5.a5}.

\bibitem{RX-Milne}
J. S. Milne, \emph{Algebraic Number Theory}, version 3.08 (2020),
Theorem~8.31, printed p.~147 (PDF index 148).
\url{https://www.jmilne.org/math/CourseNotes/ANT.pdf}.
\end{thebibliography}
